\section{How the statements were verified}
\label{app:methods}

The finite-dimensional statements of this paper fall into three kinds, and they were verified in different ways.

\paragraph{Exact statements.} Most results are identities, ranks, dimensions, spectra and decompositions of explicitly given matrices. Since every generator of the
realisation of \cref{app:realisation} is a signed permutation matrix, every blade at every rung is one as well, and products, commutators and traces are computed in integer
arithmetic without rounding. Ranks and kernel dimensions, such as the dimension $272$ of the space of mass kernels, the dimensions of centralisers and stabilisers, and the
orbit dimensions $34$ and $20$, were obtained by Gaussian elimination over the rationals. Decompositions into irreducible modules were identified by Casimir operators,
computed exactly and diagonalised on rational or algebraic eigenvalues; a statement such as $\spec C_\Delta=\{\tfrac32,\tfrac{11}2,\tfrac{15}2\}$ with multiplicities
$160,288,64$ is the exact characteristic polynomial of an integer matrix. Polynomial identities, such as the phase family of the curvature operator, the Casimir identity of the
diagonal sector, the closed form of the drift of the frame, or the Einstein tensor of the curved model on a Friedmann base, were checked as identities in the symbolic
variables, modulo the stated relations.

The exact statements were obtained by symbolic and exact computation in the Wolfram Language, and the central ones were recomputed by a second,
independent instrument: a program in pure Python using exact rational arithmetic (the \texttt{Fraction} type) and sparse representations of signed permutations, written
without reference to the first. Where possible the second instrument uses different methods: the kernel space by solving the linear conditions on all antisymmetric matrices rather than by
enumerating blades, invariant counts by enumerating weight multisets rather than by Brauer's formula, and the holomorphic quartic in an independent tensor model. Where the two
instruments use different normalisations, as for the invariant pairing in the Klein correspondence, the comparison was made on normalisation-independent quantities such as the
duality angle.

\paragraph{Derived statements.} Some statements follow from exact ones by an argument given in the text, for instance the identification of an unbroken $\uu1$ with
hypercharge, the physical reading of the second-rung vacuum, or the statement that a non-orthogonal field carries no mass. These are proved in the text and do not rest on
computation beyond the exact statements they cite.

\paragraph{Numerical statements.} A small number of statements rest on floating-point computation, and they are identified as such where they occur: the random-start
minimisations of the one-loop energy on $\mathcal T$ and on the full kernel space, which locate the competing configuration and the crossing value $x_*=24.88405$; the gradient
flows at the benchmark couplings, which illustrate the transition at $|\kappa|=7.2$; the numerical integrations of the free rotor, which reproduce the exact drift of
\cref{ex:drift} to ten digits; the integrations of a Zeitlin truncation in \cref{sec:closure}; and a few spot checks at random spinors in \cref{sec:selector}. In every case the
numerical result is either a confirmation of an exact formula or a statement explicitly labelled as numerical, and the two conjectures of \cref{sec:selection}, the global
thresholds $192/7$ on $\mathcal T$ and $x_*$ on all of $\Kc$, rest on such minimisations and are stated as conjectures.

\paragraph{Inherited statements.} Results of the established literature, such as the anomaly cancellation of the Standard Model, the Vafa--Witten theorem, Michel's theorem
on critical orbits, the Helmholtz conditions or Lanczos's identity, are cited and not recomputed; where the paper uses them, it says so.
